\documentclass[../DoAn.tex]{subfiles}
\begin{document}

Chương này tổng kết kết quả xây dựng và kiểm thử WoodCert Auction, chỉ ra các giới hạn còn tồn tại trong phạm vi đồ án, đồng thời đề xuất những hướng phát triển tiếp theo về nghiệp vụ, xác thực, tích hợp và khả năng vận hành.

\section{Kết luận}

Đồ án WoodCert Auction được thực hiện nhằm xây dựng một hệ thống hỗ trợ giao dịch sản phẩm gỗ mỹ nghệ theo quy trình khép kín, từ tạo sản phẩm, thẩm định chất lượng, tổ chức đấu giá đến thanh toán, giao nhận và giải quyết tranh chấp. Trên cơ sở khảo sát các nền tảng được trình bày ở Chương~\ref{chapter:Related_works}, đồ án tập trung vào việc kết nối quy trình thẩm định với đấu giá trực tuyến và các nghiệp vụ sau đấu giá trong cùng một hệ thống. Phạm vi này phục vụ mục tiêu học tập, nghiên cứu và kiểm chứng giải pháp kỹ thuật, chưa hướng tới thay thế đầy đủ một nền tảng thương mại có hệ sinh thái thanh toán, vận chuyển và kiểm toán độc lập.

WoodCert Auction đã hiện thực các nhóm chức năng chính cho khách truy cập, người tham gia/người mua, người bán, thẩm định viên và quản trị viên. Người bán có thể tạo sản phẩm và gửi yêu cầu thẩm định; thẩm định viên tiếp nhận hồ sơ, lập báo cáo và phát hành chứng thư; sản phẩm đạt yêu cầu được sử dụng để tạo phiên đấu giá. Người tham gia đăng ký bằng tiền cọc và đặt giá theo thời gian thực. Khi thắng phiên và phát sinh đơn hàng, người này trở thành người mua để thanh toán và tiếp tục quy trình giao nhận. Sau thanh toán, hệ thống theo dõi việc người bán xác nhận gửi hoặc bàn giao hàng, đồng thời cho phép người mua xác nhận đã nhận hàng hoặc mở tranh chấp kèm bằng chứng. Quản trị viên thực hiện các nghiệp vụ quản lý tài khoản, danh mục, doanh thu và phân xử tranh chấp trong phạm vi quyền được cấp.

Về kỹ thuật, hệ thống được tổ chức theo kiến trúc nguyên khối phân mô-đun (modular monolith), qua đó duy trì một đơn vị triển khai nhưng phân tách mã nguồn theo miền nghiệp vụ. Phiên đấu giá ở trạng thái \texttt{ACTIVE} sử dụng Redis và đoạn mã Lua để kiểm tra, cập nhật lượt đặt giá nguyên tử; MySQL lưu dữ liệu bền vững và trạng thái kết thúc. Phân hệ tài chính kết hợp khóa thao tác nghiệp vụ, ràng buộc duy nhất, giao dịch cơ sở dữ liệu và khóa lạc quan để hạn chế xử lý trùng lặp hoặc ghi đè số dư. Các tác vụ nền theo lịch tự động mở và đóng phiên, quyết toán tiền cọc, xử lý quá hạn thanh toán, quá hạn xác nhận gửi hàng và tự động hoàn tất giao nhận theo điều kiện cấu hình. Những giải pháp và đánh đổi tương ứng đã được phân tích tại Chương~\ref{chapter:SolutionAndContribution}.

Kết quả kiểm thử tại Chương~\ref{chapter:Experiment} ghi nhận 403 ca kiểm thử phía máy chủ và 196 kiểm thử đơn vị giao diện đều đạt kết quả tốt, không xảy ra lỗi thực thi. Các bước kiểm tra quy tắc mã nguồn, kiểm tra kiểu dữ liệu, đóng gói ứng dụng và 7/7 kịch bản đầu cuối bằng Playwright đều hoàn thành. Các kết quả này cung cấp bằng chứng cho những hành vi đã được kiểm thử, nhưng không thay thế kiểm toán bảo mật, kiểm toán tài chính hoặc đo kiểm tải thực tế. Hệ thống đã được triển khai production và các chức năng chính vẫn hoạt động bình thường theo xác nhận của chủ dự án; xác nhận này không được xem là số liệu định lượng về thời gian hoạt động, khả năng chịu tải hoặc cam kết dịch vụ.

Mặc dù các mục tiêu cốt lõi đã được đáp ứng trong phạm vi đồ án, WoodCert Auction vẫn còn một số giới hạn:

\noindent\textit{(i)} \textbf{Mức xác thực của chứng thư.} Mã băm SHA-256 hiện là dấu vân tay tham chiếu được lưu cùng báo cáo thẩm định. Hệ thống chưa tự động tái tạo dữ liệu đầu vào để đối chiếu khi truy vấn, chưa áp dụng chữ ký số hoặc hạ tầng khóa công khai và không bảo đảm chống chối bỏ. Độ tin cậy của kết luận chuyên môn vẫn phụ thuộc vào thẩm định viên và quy trình kiểm tra ngoài phần mềm.

\noindent\textit{(ii)} \textbf{Thanh toán, vận chuyển và trách nhiệm của người bán.} VNPay mới được tích hợp trong môi trường thử nghiệm cho chiều nạp tiền vào ví; hệ thống chưa hỗ trợ rút tiền tự động về ngân hàng hoặc đối soát hai chiều với cổng thanh toán. Trạng thái giao nhận được cập nhật thủ công, chưa kết nối đơn vị vận chuyển. Khi người bán quá hạn xác nhận gửi hoặc bàn giao hàng, hệ thống hoàn tiền cho người mua nhưng chưa có khoản ký quỹ riêng, chế tài tài chính hoặc cơ chế giảm điểm uy tín dành cho người bán.

\noindent\textit{(iii)} \textbf{Thông báo và tranh chấp.} Giao diện hiện sử dụng thông báo tức thời, chưa có trung tâm thông báo lưu bền. Hồ sơ tranh chấp được cập nhật theo chu kỳ 10 giây thay vì qua kênh thời gian thực; kết quả phân xử chỉ hỗ trợ người mua thắng hoặc người bán thắng, chưa có hoàn tiền một phần. Mã nguồn đã định nghĩa loại tệp video đóng gói, nhưng chức năng này chưa được tích hợp thành luồng tải lên và gắn bằng chứng hoàn chỉnh trong nghiệp vụ giao nhận.

\noindent\textit{(iv)} \textbf{Nhất quán và khả năng vận hành.} Sau khi Redis chấp nhận lượt đặt giá, lịch sử và bản chụp trạng thái được ghi sang MySQL theo cơ chế cố gắng ghi nhưng không bảo đảm thành công. Tác vụ nền có thể sửa bản chụp trạng thái nhưng không tái tạo được bản ghi lịch sử đã ghi thất bại. Đồ án chưa thực hiện đo kiểm tải để xác lập giới hạn mở rộng và chưa có bộ chỉ số giám sát đầy đủ cho Redis, MySQL cùng các luồng nghiệp vụ chính.

Quá trình xây dựng hệ thống cho thấy việc lựa chọn công nghệ chỉ tạo ra nền tảng ban đầu; tính đúng của nghiệp vụ còn phụ thuộc vào ranh giới dữ liệu, mô hình trạng thái, biên giao dịch và cơ chế phục hồi. Đối với luồng tài chính, tính lũy đẳng và kiểm soát cập nhật đồng thời phải được thiết kế ngay từ thao tác nghiệp vụ. Đối với đấu giá thời gian thực, việc phân tách Redis và MySQL cần đi kèm thứ tự ghi, nhật ký cảnh báo và giới hạn phục hồi rõ ràng. Đây là những kinh nghiệm chính mà sinh viên rút ra trong quá trình phân tích, triển khai và kiểm thử WoodCert Auction.

\section{Hướng phát triển}

Các hướng phát triển tiếp theo được xác định từ những giới hạn của phiên bản hiện tại và được nhóm theo mức độ liên quan nghiệp vụ.

\noindent\textit{(i)} \textbf{Hoàn thiện thông báo và trải nghiệm theo dõi.} Xây dựng trung tâm thông báo lưu bền cho các sự kiện như kết quả đấu giá, biến động số dư, hạn thanh toán và trạng thái giao nhận. Thông báo đẩy qua trình duyệt có thể được bổ sung cho sự kiện quan trọng, trong khi giao diện vẫn cho phép người dùng xem lại lịch sử đã đọc hoặc chưa đọc.

\noindent\textit{(ii)} \textbf{Mở rộng giao nhận và giải quyết tranh chấp.} Hoàn thiện luồng video đóng gói, hỗ trợ gắn bằng chứng với đơn hàng và hồ sơ tranh chấp; bổ sung phương án hoàn tiền một phần khi kết quả không phù hợp với hai lựa chọn toàn phần hiện tại. Kênh cập nhật gần thời gian thực có thể được áp dụng cho hội thoại tranh chấp nếu nhu cầu sử dụng thực tế chứng minh việc thăm dò định kỳ không còn phù hợp.

\noindent\textit{(iii)} \textbf{Hoàn thiện thanh toán và vận chuyển.} Tích hợp cổng thanh toán thực tế, quy trình rút tiền, đối soát tài chính định kỳ và cơ chế xử lý sai lệch. Việc kết nối đơn vị vận chuyển có thể hỗ trợ tạo vận đơn, theo dõi hành trình và cập nhật trạng thái tự động. Đồng thời, hệ thống cần xây dựng chính sách ký quỹ hoặc chế tài dành cho người bán, kèm điều kiện miễn trừ và quy trình khiếu nại rõ ràng.

\noindent\textit{(iv)} \textbf{Tăng cường xác thực và khả năng vận hành.} Bổ sung bước kiểm tra mã băm tự động, quản lý phiên bản dữ liệu chứng thư và nghiên cứu chữ ký số để nâng mức xác thực. Hệ thống cần thu thập chỉ số, nhật ký và cảnh báo cho các thành phần chính; thực hiện đo kiểm tải theo kịch bản tái lập được để xác định giới hạn trước khi quyết định mở rộng tài nguyên hoặc tách dịch vụ.

Các định hướng trên không thay đổi phạm vi kết quả hiện tại của đồ án. Chúng là cơ sở để WoodCert Auction tiếp tục phát triển từ một hệ thống đã hoàn thiện các luồng cốt lõi thành nền tảng có mức tự động hóa, khả năng kiểm chứng và năng lực vận hành cao hơn.

\end{document}
